
\documentclass{article}
\usepackage{colortbl}
\usepackage{makecell}
\usepackage{multirow}
\usepackage{supertabular}

\begin{document}

\newcounter{utterance}

\centering \large Interaction Transcript for game `codenames', experiment `frequency\_high', episode 6 with reflexion{-}Qwen3.5{-}27B.
\vspace{24pt}

{ \footnotesize  \setcounter{utterance}{1}
\setlength{\tabcolsep}{0pt}
\begin{supertabular}{c@{$\;$}|p{.15\linewidth}@{}p{.15\linewidth}p{.15\linewidth}p{.15\linewidth}p{.15\linewidth}p{.15\linewidth}}
    \# & $\;$A & \multicolumn{4}{c}{Game Master} & $\;\:$B\\
    \hline

    \theutterance \stepcounter{utterance}  
    & & & \multicolumn{2}{p{0.3\linewidth}}{
        \cellcolor[rgb]{0.9,0.9,0.9}{
            \makecell[{{p{\linewidth}}}]{
                \texttt{\tiny{[GM$|$GM]}}
                {'hidden': {'team': ['horse', 'mama', 'job', 'situation', 'key', 'pain', 'road', 'attention', 'bar'], 'innocent': ['school', 'doctor', 'world', 'lord', 'uncle', 'bed', 'stuff'], 'opponent': ['water', 'girl', 'couple', 'floor', 'evidence', 'club', 'honey', 'death'], 'assassin': ['window']}, 'revealed': {'team': {'team': [], 'innocent': [], 'opponent': [], 'assassin': []}, 'opponent': {'team': [], 'innocent': [], 'opponent': [], 'assassin': []}}}
            }
        }
    }
    & & \\ \\

    \theutterance \stepcounter{utterance}  
    & & \multicolumn{4}{p{0.6\linewidth}}{
        \cellcolor[rgb]{0.9,0.9,0.9}{
            \makecell[{{p{\linewidth}}}]{
                \texttt{\tiny{[P1$\langle$GM]}}
                \texttt{You have attempted this same task before and did not succeed. These are the plans you wrote after those attempts.} \\
\texttt{Trial \#1: New plan: The previous attempt was aborted because the clue "key" was identical to one of the target words ("key"), violating the rule that the clue cannot be one of the words in the lists. Additionally, the clue must not contain parts of the target words. Next time, I will select a clue that is semantically related to multiple target words (e.g., "horse", "mama", "job", "situation", "key", "pain", "road", "attention", "bar") but is distinct from all of them and does not contain any of their substrings. For example, I could use "life" or "problem" as a clue, ensuring the TARGETS list only includes words from my team's list that fit the clue, while strictly avoiding the kill word "window" and any opponent/distractor words. I will double{-}check that the clue word is not present in the target list before submitting.} \\
\\ 
\texttt{Trial \#2: New plan: The previous attempt was aborted because the TARGETS list included "attention" and "bar," which are weak or invalid associations for the clue "LIFE" in the context of this specific game's strict semantic rules, or potentially the system flagged the inclusion of words that do not strongly fit the clue as a format violation regarding the "related" constraint. More critically, I must ensure the TARGETS list contains *only* words that have a direct, unambiguous semantic link to the clue to avoid rejection. I will select a tighter clue like "WORK" which strongly relates to "job" and "bar" (as in a workplace or place of work) and "horse" (workhorse), or "PAIN" which relates to "pain" (wait, cannot use the word itself), so "PAIN" is bad. Let's try "CARE". "Mama" (caregiver), "job" (care for work), "attention" (care), "horse" (care for animals). Actually, the safest bet is to pick a clue that covers a smaller, very tight subset to ensure validity. I will choose the clue "WORK" and list only "job" and "horse" (as in workhorse) to ensure high precision, or "MAMA" is a target, so clue "MOM" is bad (contains part). Let's try "ROLE". Targets: "job", "situation", "mama" (role). No, "mama" is a role. Better: Clue "ISSUE". Targets: "situation", "pain", "problem" (not in list). Let's go with "CARE". Targets: "mama", "attention", "horse". I will strictly verify that the clue is not in the target list, does not contain substrings of targets, and that every listed target is a strong semantic match. I will output the format exactly as requested.} \\
\\ 
\texttt{Revised specific plan: I will use the clue "CARE". It relates to "mama" (motherly care), "attention" (paying care), and "horse" (animal care). I will list only these three to ensure the association is undeniable. I will double{-}check that "CARE" is not in the target list (it isn't) and does not contain parts of "horse", "mama", "job", "situation", "key", "pain", "road", "attention", "bar". "Care" is not a substring of any. I will format the output strictly as:} \\
\texttt{CLUE: CARE} \\
\texttt{TARGETS: mama, attention, horse} \\
\\ 
\\ 
\texttt{Provide one single{-}word clue that relates to some of your team's words from the list below. You can choose to how many and to which words your clue relates to.} \\
\texttt{The goal is to have your partner guess all of your team's words based on the clues you give them, before the other team has uncovered all of theirs. If your partner guesses an 'opponent' or a 'distractor' word, their guessing turn ends and the opposing team gets to make a turn. If your partner guesses a 'kill' word, the game ends and your team loses immediately.} \\
\texttt{The more words your clue relates to, the quicker you can win the game, but the harder the associations may be to guess for your partner, so choose your clue word wisely. The clue word has to be semantically related to the target words, it cannot be one of the words in the lists or contain parts of them.} \\
\\ 
\texttt{Always give your single{-}word clue and your comma{-}separated list of related target words in the following format and make your answers as short as possible, never include any other text than is required in this form:} \\
\\ 
\texttt{CLUE: <WORD>} \\
\texttt{TARGETS: <TARGETS>} \\
\\ 
\texttt{Your team words are: horse, mama, job, situation, key, pain, road, attention, bar.} \\
\texttt{Your opponent's team words are: water, girl, couple, floor, evidence, club, honey, death.} \\
\texttt{Distractor words are: school, doctor, world, lord, uncle, bed, stuff.} \\
\texttt{The kill words are: window.} \\
            }
        }
    }
    & \\ \\

    \theutterance \stepcounter{utterance}  
    & \multicolumn{4}{p{0.6\linewidth}}{
        \cellcolor[rgb]{0.9,0.9,0.9}{
            \makecell[{{p{\linewidth}}}]{
                \texttt{\tiny{[P1$\rangle$GM]}}
                \texttt{CLUE: PAIN} \\
\texttt{TARGETS: pain, situation, job, road, horse, bar, attention, key, mama} \\
            }
        }
    }
    & & \\ \\

    \theutterance \stepcounter{utterance}  
    & & & \multicolumn{2}{p{0.3\linewidth}}{
        \cellcolor[rgb]{0.9,0.9,0.9}{
            \makecell[{{p{\linewidth}}}]{
                \texttt{\tiny{[GM$|$GM]}}
                {'player': 'Cluegiver', 'type': 'clue is morphologically related to word on the board', 'utterance': 'CLUE: PAIN\nTARGETS: pain, situation, job, road, horse, bar, attention, key, mama', 'message': "Your clue 'pain' is morphologically similar to the word pain, please choose another clue word.", 'clue': 'pain', 'similar_board_word': 'pain'}
            }
        }
    }
    & & \\ \\

    \theutterance \stepcounter{utterance}  
    & \multicolumn{4}{p{0.6\linewidth}}{
        \cellcolor[rgb]{0.95,0.95,0.95}{
            \makecell[{{p{\linewidth}}}]{
                \texttt{\tiny{[P2$\rightarrow$P1]}}
                \texttt{Post{-}episode Reflexion (not sent to GM):} \\
\texttt{New plan: The previous attempt was aborted because the clue "PAIN" was identical to one of the target words ("pain"), violating the explicit rule that the clue cannot be one of the words in the lists. Additionally, listing all nine words for a single clue is risky if the semantic link isn't perfect for every single one, potentially triggering a validation error for weak associations. Next time, I will select a clue that is semantically related to a subset of my team's words but is strictly distinct from all words in the game (targets, opponents, distractors, and kill words) and does not contain any of them as substrings. I will choose the clue "WORK" which strongly relates to "job" and "horse" (as in workhorse) and "bar" (as in a place of work or a bar of soap/metal, though "job" and "horse" are the strongest). To be safe and precise, I will limit the TARGETS list to only the words with the strongest, unambiguous connection to the clue. I will verify the clue is not "window" (kill word) and not in the opponent/distractor lists. My response will be strictly formatted as:} \\
\texttt{CLUE: WORK} \\
\texttt{TARGETS: job, horse} \\
            }
        }
    }
    & & \\ \\

\end{supertabular}
}

\end{document}
